% ============================================================
%  Capítulo 1 — Introducción a la Inferencia Estadística
% ============================================================

\section{Tópicos en Probabilidad}

Para comprender la inferencia estadística, es fundamental contar con una base sólida en probabilidad. En esta sección se introducen los conceptos que establecerán el marco teórico del curso.

\subsection{Espacios de trabajo}

Sea $(\Omega,\mathcal{F},P)$ un \textbf{espacio de probabilidad}, donde:
\begin{itemize}
\item $\Omega$ es el espacio muestral,
\item $\mathcal{F}$ es una $\sigma$-álgebra de subconjuntos de $\Omega$,
\item $P$ es una medida de probabilidad.
\end{itemize}

Consideraremos además un \textbf{espacio medible} $(\mathcal{E},\mathcal{B}(\mathcal{E}))$, donde $\mathcal{E}$ es el espacio de valores y $\mathcal{B}(\mathcal{E})$ su $\sigma$-álgebra de Borel.

\subsection{Variable aleatoria y ley}

\begin{defi}[Variable aleatoria]
Una variable aleatoria con valores en $\mathcal{E}$ es una función medible
\[
X:\Omega \to \mathcal{E}.
\]
\end{defi}

La medibilidad significa que, para todo $A \in \mathcal{B}(\mathcal{E})$,
\[
X^{-1}(A) = \{\omega \in \Omega : X(\omega) \in A\} \in \mathcal{F}.
\]

\begin{defi}[Ley de una variable aleatoria]
La \textbf{ley} de $X$ es la medida de probabilidad $P_X$ sobre $\mathcal{E}$ definida por
\[
P_X(A) = P(X^{-1}(A)), \quad \forall A \in \mathcal{B}(\mathcal{E}).
\]
\end{defi}

Equivalente y compactamente, se escribe
\[
P_X = P \circ X^{-1},
\]
y se conoce como la \textbf{medida imagen} (o \textit{pushforward}) de $P$ por $X$.

\subsection{Función de distribución acumulada}

Cuando $\mathcal{E} = \mathbb{R}$, la ley de $X$ puede describirse mediante su función de distribución acumulada.

\begin{defi}[Función de distribución acumulada (FDA)]
La FDA de $X$ es la función $F_X:\mathbb{R} \to [0,1]$ definida por
\[
\begin{aligned}
F_X : \mathbb{R} &\to [0,1] \\
x &\mapsto P(X \leq x) = P_X((-\infty, x]).
\end{aligned}
\]
\end{defi}

La FDA determina completamente la ley de $X$, ya que permite reconstruir la medida $P_X$ sobre la $\sigma$-álgebra de Borel.

\subsection{Esperanza}

\begin{defi}[Esperanza]
Si $X$ es integrable, su esperanza se define por
\[
\E[X] = \int_{\Omega} X(\omega) \, dP(\omega).
\]
\end{defi}

Mediante el Teorema del Transporte, esta puede expresarse en términos de la ley de $X$:
\[
\E[X] = \int_{\mathbb{R}} x \, dP_X(x).
\]

\subsection{Espacios \texorpdfstring{$L^p$}{Lp}}

Surge naturalmente la pregunta: \textit{¿para qué variables aleatorias está bien definida la esperanza?}

Para responderla, introducimos los espacios $L^p$. Para $1 \leq p < \infty$, definimos
\[
L^p(\Omega,\mathcal{F},P)
= \left\{ X : \E[|X|^p] < \infty \right\}.
\]

En particular:
\begin{itemize}
\item $X \in L^1$ si y solo si $\E[|X|] < \infty$,
\item en este caso, la esperanza $\E[X]$ está bien definida.
\end{itemize}

Además, como $P$ es una medida de probabilidad (finita), se tiene que
\[
L^p \subseteq L^1 \quad \text{para todo } p \geq 1.
\]

Sin embargo, la inclusión es estricta: existen variables en $L^1$ que no pertenecen a ningún $L^p$ con $p > 1$.

\medskip

\noindent
\textbf{Conclusión:} el espacio natural donde la esperanza está bien definida es $L^1(\Omega,\mathcal{F},P)$.

\subsection{Covarianza y Varianza}

Primero las definimos:

\begin{defi}[Covarianza]
La covarianza entre dos variables aleatorias $X$ e $Y$ es definida por
\[  \Cov(X,Y) = \E[(X - \E[X])(Y - \E[Y])]. \]
\end{defi}    

\begin{defi}[Varianza]
La varianza de una variable aleatoria $X$ es definida por
\[  \Var(X) = \E[(X - \E[X])^2]. \]
\end{defi}    

\textbf{Observación:} la varianza es un caso particular de covarianza, ya que $\Var(X) = \Cov(X,X)$ y además si $X, Y \in L^2$ entonces $\langle X,Y\rangle_{L^2} = \Cov(X,Y) = \E[XY] - \E[X]\E[Y]$.

\subsection{Leyes clásicas:}

\textbf{Distribuciones discretas} ($\mathcal{E} = \mathbb{R}$):

\begin{itemize}
\item \textbf{Bernoulli:} $X \sim \text{Bern}(p)$ con $P(X=1) = p$ y $P(X=0) = 1-p$. La ley es:
\[
P_X = p\,\delta_1 + (1-p)\,\delta_0,
\]
y por ende, la esperanza:
\[
\E[X] = \int_{\mathbb{R}} x \, dP_X(x) = \int_{\mathbb{R}} x \, d(p\,\delta_1 + (1-p)\,\delta_0)(x) = p.
\]

\item \textbf{Binomial:} $X \sim \text{Bin}(n,p)$ con $n \in \mathbb{N}$ y $p \in [0,1]$. La ley es:
\[
P_X = \sum_{k=0}^{n} \binom{n}{k} p^k (1-p)^{n-k} \, \delta_k,
\]
y por ende, la esperanza:
\[
\E[X] = \sum_{k=0}^{n} k \, dP_X(k) = \sum_{k=0}^{n} k \binom{n}{k} p^k (1-p)^{n-k} = np.
\]

\item \textbf{Poisson:} $X \sim \text{Poi}(\lambda)$ con $\lambda > 0$. La ley es:
\[
P_X = \sum_{k=0}^{\infty} e^{-\lambda} \frac{\lambda^k}{k!} \, \delta_k,
\]
y por ende, la esperanza:
\[
\E[X] = \sum_{k=0}^{\infty} k \, dP_X(k) = \sum_{k=0}^{\infty} k \, e^{-\lambda} \frac{\lambda^k}{k!} = \lambda.
\]
\end{itemize}

\subsection{Leyes absolutamente continuas}

Hasta ahora hemos descrito leyes discretas mediante combinaciones de deltas de Dirac. Sin embargo, muchas variables aleatorias de interés toman valores en $\mathbb{R}$ de manera continua.

Sea $\lambda$ la \textbf{medida de Lebesgue} en $\mathbb{R}$.

\begin{defi}[Medida absolutamente continua]
Decimos que la ley $P_X$ es \textbf{absolutamente continua} respecto de la medida de Lebesgue, y escribimos
\[
P_X \ll \lambda,
\]
si para todo conjunto boreliano $A$,
\[
\lambda(A) = 0 \implies P_X(A) = 0.
\]
\end{defi}

Este concepto es fundamental, ya que permite introducir densidades.

\begin{teo}[Radon--Nikodym]
Si $P_X \ll \lambda$, entonces existe una función medible $f_X:\mathbb{R} \to [0,\infty)$ tal que
\[
P_X(A) = \int_A f_X(x)\,dx, \quad \forall A \in \mathcal{B}(\mathbb{R}).
\]
\end{teo}

La función $f_X$ se denomina \textbf{densidad} de $X$, y se escribe formalmente
\[
P_X(dx) = f_X(x)\,dx.
\]

\subsubsection*{Consecuencias}

Si $X$ admite densidad $f_X$, entonces:

\begin{itemize}
\item La probabilidad de un conjunto se calcula como
\[
P(X \in A) = \int_A f_X(x)\,dx.
\]

\item La esperanza puede escribirse como
\[
\E[X] = \int_{\mathbb{R}} x f_X(x)\,dx.
\]
\end{itemize}

\subsubsection*{Ejemplos clásicos}

\begin{itemize}
\item \textbf{Uniforme:} $X \sim \mathcal{U}(a,b)$
\[
f_X(x) = \frac{1}{b-a}\,\mathbb{1}_{[a,b]}(x)
\]

\item \textbf{Exponencial:} $X \sim \mathrm{Exp}(\lambda)$
\[
f_X(x) = \lambda e^{-\lambda x}\,\mathbb{1}_{[0,\infty)}(x)
\]

\item \textbf{Normal:} $X \sim \mathcal{N}(\mu,\sigma^2)$
\[
f_X(x) = \frac{1}{\sqrt{2\pi\sigma^2}}
\exp\left(-\frac{(x-\mu)^2}{2\sigma^2}\right)
\]
\end{itemize}

\medskip

\noindent
\textbf{Interpretación:} en el caso discreto, la ley se escribe como suma de deltas; en el caso continuo, como una densidad respecto a la medida de Lebesgue. Ambos enfoques son unificados mediante la noción de medida.

\subsection{Independencia}

La noción de independencia puede formularse tanto en términos de eventos como en términos de leyes.

\medskip

\noindent
\textbf{Recordatorio:} si $X:\Omega \to \mathcal{E}$ es una variable aleatoria, entonces $X$ induce la $\sigma$-álgebra
\[
\sigma(X)=\{X^{-1}(A):A\in\mathcal{B}(\mathcal{E})\}.
\]
En este sentido, la independencia de variables aleatorias puede entenderse como independencia entre las $\sigma$-álgebras que ellas generan.

\begin{defi}[Variables aleatorias independientes]
Sean $X_1,\dots,X_n$ variables aleatorias con valores en $\mathcal{E}$. Diremos que son \textbf{independientes} si para todo $B_1,\dots,B_n \in \mathcal{B}(\mathcal{E})$ se tiene
\[
P(X_1\in B_1,\dots,X_n\in B_n)
=
\prod_{i=1}^{n} P(X_i\in B_i).
\]
\end{defi}

\noindent
Dicho de otra forma, la ley de la variable aleatoria vectorial
\[
(X_1,\dots,X_n):\Omega\to \mathcal{E}^n
\]
satisface
\[
P_{(X_1,\dots,X_n)}
=
P_{X_1}\otimes P_{X_2}\otimes\cdots\otimes P_{X_n}.
\]

En efecto, para todo $B_1,\dots,B_n\in\mathcal{B}(\mathcal{E})$,
\[
P_{(X_1,\dots,X_n)}(B_1\times\cdots\times B_n)
=
P\big((X_1,\dots,X_n)\in B_1\times\cdots\times B_n\big),
\]
y por definición de la variable aleatoria vectorial,
\[
P\big((X_1,\dots,X_n)\in B_1\times\cdots\times B_n\big)
=
P(X_1\in B_1,\dots,X_n\in B_n).
\]

Por otra parte,
\[
(P_{X_1}\otimes\cdots\otimes P_{X_n})(B_1\times\cdots\times B_n)
=
\prod_{i=1}^n P_{X_i}(B_i)
=
\prod_{i=1}^n P(X_i\in B_i).
\]

Luego, la independencia equivale a que la ley conjunta coincida con el producto de las leyes marginales.

\medskip

\noindent
\textbf{Observación:} no debe confundirse independencia conjunta con independencia dos a dos.

Decimos que $X_1,\dots,X_n$ son \textbf{independientes dos a dos} si para todo $i\neq j$, las variables $X_i$ y $X_j$ son independientes.

En general,
\[
\text{independencia conjunta} \implies \text{independencia dos a dos},
\]
pero el recíproco es falso.

\subsubsection*{Ejemplo}

Sean $X$ e $Y$ variables aleatorias i.i.d.\ con ley $\mathrm{Bernoulli}(1/2)$ y definamos
\[
Z = X+Y \pmod 2.
\]

Entonces $(X,Y,Z)$ son independientes dos a dos, pero no son independientes conjuntamente.

En efecto, existe una relación determinista entre ellas:
\[
Z = X+Y \pmod 2.
\]
Por tanto, conocer $X$ e $Y$ determina completamente a $Z$, lo que impide la independencia mutua.

Por ejemplo,
\[
P(X=0,Y=0,Z=1)=0,
\]
mientras que
\[
P(X=0)\,P(Y=0)\,P(Z=1)
=
\frac12\cdot\frac12\cdot\frac12
=
\frac18.
\]
Como
\[
0\neq \frac18,
\]
se concluye que $X$, $Y$ y $Z$ no son independientes.

\subsubsection*{Factorización de esperanzas}

La independencia permite factorizar esperanzas de productos de funciones, como veremos a continuación:

\textbf{Proposición:} si $X_1,\dots,X_n$ son independientes y $h_1,\dots,h_n:\mathcal{E}\to\mathbb{R}$ son medibles y no negativas, entonces
\[
\E\left[\prod_{i=1}^n h_i(X_i)\right]
=
\prod_{i=1}^n \E[h_i(X_i)].
\]

Más generalmente, la igualdad sigue siendo válida siempre que las cantidades involucradas estén bien definidas.

\medskip

\noindent
\subsubsection*{Demostración}

Sea
\[
H:\mathcal{E}^n \to \mathbb{R}
\]
definida por
\[
H(x_1,\dots,x_n)=h_1(x_1)\,h_2(x_2)\cdots h_n(x_n).
\]

Luego,
\[
\E\left[\prod_{i=1}^n h_i(X_i)\right]
=
\E\big(H(X_1,\dots,X_n)\big).
\]

Por el Teorema del Transporte,
\[
\E\big(H(X_1,\dots,X_n)\big)
=
\int_{\mathcal{E}^n} H(x_1,\dots,x_n)\,dP_{(X_1,\dots,X_n)}(x_1,\dots,x_n).
\]

Como $X_1,\dots,X_n$ son independientes,
\[
P_{(X_1,\dots,X_n)}
=
P_{X_1}\otimes\cdots\otimes P_{X_n}.
\]

Por lo tanto,
\[
\int_{\mathcal{E}^n} H(x_1,\dots,x_n)\,dP_{(X_1,\dots,X_n)}
=
\int_{\mathcal{E}^n} h_1(x_1)\cdots h_n(x_n)\,d(P_{X_1}\otimes\cdots\otimes P_{X_n}).
\]

Aplicando Tonelli--Fubini,
\[
\int_{\mathcal{E}^n} h_1(x_1)\cdots h_n(x_n)\,d(P_{X_1}\otimes\cdots\otimes P_{X_n})
\]
\[
=
\int_{\mathcal{E}} h_1(x_1)\,dP_{X_1}(x_1)\cdots
\int_{\mathcal{E}} h_n(x_n)\,dP_{X_n}(x_n).
\]

Finalmente, para cada $i$,
\[
\int_{\mathcal{E}} h_i(x_i)\,dP_{X_i}(x_i)=\E[h_i(X_i)].
\]

Así,
\[
\E\left[\prod_{i=1}^n h_i(X_i)\right]
=
\prod_{i=1}^n \E[h_i(X_i)].
\]
\hfill $\square$

\medskip

\noindent
\textbf{Observación:} el resultado sigue siendo válido si
\[
\E[|h_i(X_i)|]<\infty,
\qquad i=1,\dots,n.
\]
En ese caso,
\[
\E\left[\prod_{i=1}^n |h_i(X_i)|\right]
=
\prod_{i=1}^n \E[|h_i(X_i)|]
<\infty,
\]
por lo que el producto es integrable y puede aplicarse Fubini.

\medskip

\noindent
\textbf{Corolario:} si $X\perp Y$ y $X,Y\in L^2$, entonces
\[
\Cov(X,Y)=0.
\]

\textbf{Demostración:} por independencia,
\[
\E[XY]=\E[X]\E[Y].
\]
Luego,
\[
\Cov(X,Y)=\E[XY]-\E[X]\E[Y]=0.
\]
\hfill $\square$

\medskip

\noindent
\textbf{Última observación:} la recíproca es falsa en general:
\[
\Cov(X,Y)=0 \not\implies X\perp Y.
\]

\subsubsection*{Recordatorio}

Recordemos que, si $X:\Omega\to\mathbb{R}$ es una variable aleatoria y $h:\mathbb{R}\to\mathbb{R}$ es medible, entonces
\[
\E[h(X)] = \int_{\mathbb{R}} h(x)\,P_X(dx),
\]
siempre que la integral esté bien definida.

\medskip

\noindent
\textbf{(1) Caso absolutamente continuo.}  
Si $X$ es absolutamente continua, entonces existe una densidad $f_X$ tal que
\[
P_X(dx)=f_X(x)\,dx.
\]
Luego,
\[
\E[h(X)]
=
\int_{\mathbb{R}} h(x)\,P_X(dx)
=
\int_{\mathbb{R}} h(x)f_X(x)\,dx.
\]

\medskip

\noindent
\textbf{(2) Caso discreto.}  
Si la ley de $X$ es discreta, por ejemplo con soporte en $\mathbb{N}$, entonces
\[
P_X(dx)=\sum_{k\in\mathbb{N}} P(X=k)\,\delta_k(dx).
\]
Por lo tanto,
\[
\E[h(X)]
=
\int_{\mathbb{R}} h(x)\,P_X(dx)
\]
\[
=
\int_{\mathbb{R}} h(x)\left(\sum_{k\in\mathbb{N}} P(X=k)\,\delta_k(dx)\right).
\]
Intercambiando suma e integral,
\[
\E[h(X)]
=
\sum_{k\in\mathbb{N}} P(X=k)\int_{\mathbb{R}} h(x)\,\delta_k(dx).
\]
Como
\[
\int_{\mathbb{R}} h(x)\,\delta_k(dx)=h(k),
\]
se obtiene
\[
\E[h(X)]
=
\sum_{k\in\mathbb{N}} h(k)\,P(X=k).
\]

\subsection{Función característica}

Sea $X$ una variable aleatoria con valores en $\mathbb{R}^d$, es decir,
\[
X=(X_1,\dots,X_d).
\]

La \textbf{función característica} de $X$ es la función
\[
\varphi_X:\mathbb{R}^d\to\mathbb{C}
\]
definida por
\[
\varphi_X(t)=\E\big(e^{i\langle t,X\rangle}\big),
\qquad t\in\mathbb{R}^d.
\]

Si $t=(t_1,\dots,t_d)$, entonces
\[
\langle t,X\rangle=\sum_{j=1}^d t_jX_j,
\]
y por lo tanto
\[
\varphi_X(t)
=
\E\left(\exp\left(i\sum_{j=1}^d t_jX_j\right)\right).
\]

\medskip

\noindent
\textbf{Observación:} la función característica puede verse como una versión probabilística de la transformada de Fourier de la ley de $X$, pues
\[
\varphi_X(t)=\int_{\mathbb{R}^d} e^{i\langle t,x\rangle}\,P_X(dx).
\]

\subsubsection*{Casos particulares}

\textbf{Si $X$ tiene densidad $f_X$ en $\mathbb{R}^d$, entonces}
\[
\varphi_X(t)
=
\int_{\mathbb{R}^d} e^{i\langle t,x\rangle}f_X(x)\,dx.
\]

\textbf{Si $X$ es discreta con valores $x_k\in\mathbb{R}^d$, entonces}
\[
\varphi_X(t)
=
\sum_k e^{i\langle t,x_k\rangle}P(X=x_k).
\]

\subsubsection*{Observación}

Si $X$ es una variable aleatoria real, entonces su función característica viene dada por
\[
\varphi_X(t)=\E(e^{itX})=\int_{\mathbb{R}} e^{itx}\,P_X(dx).
\]

Si además $X$ es absolutamente continua con densidad $f_X$, entonces
\[
P_X(dx)=f_X(x)\,dx,
\]
y por lo tanto
\[
\varphi_X(t)
=
\int_{\mathbb{R}} e^{itx}\,P_X(dx)
=
\int_{\mathbb{R}} e^{itx}f_X(x)\,dx.
\]

\subsubsection*{Ejemplos}

\textbf{(1) Si $X\sim \mathrm{Bernoulli}(p)$, entonces}
\[
\varphi_X(t)=\E(e^{itX}).
\]
Como $X$ solo puede tomar los valores $0$ y $1$, se tiene
\[
\varphi_X(t)
=
e^{it}\,P(X=1)+e^{i0t}\,P(X=0)
=
pe^{it}+(1-p).
\]
Es decir,
\[
\varphi_X(t)=(1-p)+pe^{it}.
\]

\medskip

\textbf{(2) Si $Y\sim \mathrm{Binomial}(N,p)$, entonces}

Recordemos que una variable binomial puede escribirse como suma de variables Bernoulli i.i.d.:
\[
Y=\sum_{j=1}^{N} Z_j,
\qquad Z_j \stackrel{\mathrm{i.i.d.}}{\sim} \mathrm{Bernoulli}(p).
\]

Luego,
\[
\varphi_Y(t)=\E(e^{itY})
=
\E\left(e^{it\sum_{j=1}^{N} Z_j}\right)
=
\E\left(\prod_{j=1}^{N} e^{itZ_j}\right).
\]

Como las variables $Z_1,\dots,Z_N$ son independientes, se obtiene
\[
\varphi_Y(t)
=
\prod_{j=1}^{N}\E(e^{itZ_j})
=
\prod_{j=1}^{N}\varphi_{Z_j}(t).
\]

Como todas tienen la misma ley Bernoulli$(p)$,
\[
\varphi_Y(t)=\big(\varphi_{Z_1}(t)\big)^N
=
\big((1-p)+pe^{it}\big)^N.
\]

Por tanto,
\[
\varphi_Y(t)=\big((1-p)+pe^{it}\big)^N.
\]

\medskip

\textbf{(3) Si $Z\sim \mathcal{N}(0,1)$, queremos calcular $\varphi_Z(t)$.}

Más generalmente, si
\[
W\sim \mathcal{N}(m,\sigma^2),
\]
entonces puede escribirse como
\[
W=m+\sigma Z,
\qquad Z\sim \mathcal{N}(0,1).
\]

Así,
\[
\varphi_W(t)=\E(e^{itW})
=
\E\big(e^{it(m+\sigma Z)}\big)
=
e^{imt}\E(e^{i\sigma t Z}).
\]

Por definición de función característica,
\[
\E(e^{i\sigma t Z})=\varphi_Z(\sigma t).
\]

Luego,
\[
\varphi_W(t)=e^{imt}\varphi_Z(\sigma t).
\]

\noindent
Ahora bien, si $Z\sim \mathcal{N}(0,1)$, entonces
\[
\varphi_Z(t)=e^{-t^2/2}.
\]

Por consiguiente, si $W\sim \mathcal{N}(m,\sigma^2)$, se obtiene
\[
\varphi_W(t)=e^{imt}\varphi_Z(\sigma t)
=
e^{imt}e^{-\sigma^2 t^2/2}.
\]

Es decir,
\[
\varphi_W(t)=\exp\left(imt-\frac{\sigma^2 t^2}{2}\right).
\]

\subsection{Modos de convergencia}

Hasta ahora hemos estudiado variables aleatorias, sus leyes, independencia y función característica. El siguiente paso natural es comparar sucesiones de variables aleatorias y precisar en qué sentido una sucesión puede \emph{converger}.

Sean $X,X_1,X_2,\dots$ variables aleatorias reales definidas sobre un mismo espacio de probabilidad.

\begin{defi}[Convergencia casi segura]
Decimos que $X_n$ converge \textbf{casi seguramente} a $X$, y escribimos
\[
X_n \xrightarrow{\text{c.s.}} X,
\]
si
\[
P\big(\{\omega\in\Omega : X_n(\omega)\to X(\omega)\}\big)=1.
\]
\end{defi}

\begin{defi}[Convergencia en probabilidad]
Decimos que $X_n$ converge \textbf{en probabilidad} a $X$, y escribimos
\[
X_n \xrightarrow{P} X,
\]
si para todo $\varepsilon>0$,
\[
P(|X_n-X|>\varepsilon)\to 0.
\]
\end{defi}

\begin{defi}[Convergencia en \texorpdfstring{$L^p$}{Lp}]
Para $1\le p<\infty$, decimos que $X_n$ converge a $X$ en $L^p$, y escribimos
\[
X_n \xrightarrow{L^p} X,
\]
si
\[
\E[|X_n-X|^p]\to 0.
\]
\end{defi}

\begin{defi}[Convergencia en distribución]
Decimos que $X_n$ converge \textbf{en distribución} a $X$, y escribimos
\[
X_n \xrightarrow{d} X,
\]
si para todo punto de continuidad $x$ de la FDA $F_X$,
\[
F_{X_n}(x)\to F_X(x).
\]
\end{defi}

\subsubsection*{Relaciones entre los distintos modos de convergencia}

\textbf{Proposición:}
\[
X_n \xrightarrow{\mathrm{c.s.}} X
\quad\Longrightarrow\quad
X_n \xrightarrow{P} X
\quad\Longrightarrow\quad
X_n \xrightarrow{d} X.
\]

Además, para todo $p\ge 1$,
\[
X_n \xrightarrow{L^p} X
\quad\Longrightarrow\quad
X_n \xrightarrow{P} X.
\]

\medskip

\noindent
\textbf{Observación:} la recíproca de estas implicaciones es falsa en general.

\medskip

\noindent
\textbf{Observación importante:} si $c\in\mathbb{R}$ es una constante, entonces
\[
X_n \xrightarrow{d} c
\quad\Longleftrightarrow\quad
X_n \xrightarrow{P} c.
\]

Es decir, cuando el límite es determinista, convergencia en distribución y convergencia en probabilidad coinciden.

\subsubsection*{Criterio en términos de funciones características}

La función característica permite estudiar convergencia en distribución de manera muy eficiente.

\begin{teo}[Teorema de continuidad de Lévy]
Sean $X_n$ y $X$ variables aleatorias con valores en $\mathbb{R}^d$. Entonces
\[
X_n \xrightarrow{d} X
\]
si y solo si
\[
\varphi_{X_n}(t)\to \varphi_X(t),
\qquad \forall t\in\mathbb{R}^d.
\]
\end{teo}

Este resultado será la herramienta fundamental para demostrar el Teorema Central del Límite.

\subsection{Leyes de los grandes números}

Sea $(X_n)_{n\ge 1}$ una sucesión de variables aleatorias i.i.d.\ con
\[
\mu=\E[X_1].
\]
Definimos el promedio muestral por
\[
\overline{X}_n=\frac{1}{n}\sum_{i=1}^n X_i.
\]

La intuición básica es que, al promediar muchas observaciones independientes con la misma ley, el promedio muestral debería acercarse al valor esperado.

\subsubsection*{Ley débil de los grandes números}

\begin{teo}[Ley Débil de los Grandes Números]
Si $X_1,X_2,\dots$ son i.i.d.\ y
\[
\E[X_1]=\mu,
\qquad
\Var(X_1)=\sigma^2<\infty,
\]
entonces
\[
\overline{X}_n \xrightarrow{P} \mu.
\]
\end{teo}

\subsubsection*{Demostración}

Como las variables son i.i.d., se tiene
\[
\E[\overline{X}_n]
=
\E\left[\frac{1}{n}\sum_{i=1}^n X_i\right]
=
\frac{1}{n}\sum_{i=1}^n \E[X_i]
=
\mu.
\]

Además, por independencia,
\[
\Var(\overline{X}_n)
=
\Var\left(\frac{1}{n}\sum_{i=1}^n X_i\right)
=
\frac{1}{n^2}\sum_{i=1}^n \Var(X_i)
=
\frac{n\sigma^2}{n^2}
=
\frac{\sigma^2}{n}.
\]

Luego, por la desigualdad de Chebyshev, para todo $\varepsilon>0$,
\[
P(|\overline{X}_n-\mu|>\varepsilon)
\le
\frac{\Var(\overline{X}_n)}{\varepsilon^2}
=
\frac{\sigma^2}{n\varepsilon^2}
\to 0.
\]

Por lo tanto,
\[
\overline{X}_n \xrightarrow{P} \mu.
\]
\hfill $\square$

\subsubsection*{Ley fuerte de los grandes números}

\begin{teo}[Ley Fuerte de los Grandes Números]
Si $X_1,X_2,\dots$ son i.i.d.\ y
\[
\E[|X_1|]<\infty,
\]
entonces
\[
\overline{X}_n \xrightarrow{\mathrm{c.s.}} \mu,
\qquad \mu=\E[X_1].
\]
\end{teo}

\medskip

\noindent
\textbf{Observación:} la ley fuerte es más fuerte que la ley débil, pues la convergencia casi segura implica convergencia en probabilidad.

\subsubsection*{Interpretación estadística}

Si $X_1,\dots,X_n$ representan observaciones de una misma población, entonces la ley de los grandes números justifica que el promedio muestral
\[
\overline{X}_n
\]
sea un buen estimador de la media poblacional
\[
\mu=\E[X_1].
\]

Por ejemplo, si $X_i\sim \mathrm{Bernoulli}(p)$ i.i.d., entonces
\[
\overline{X}_n=\frac{1}{n}\sum_{i=1}^n X_i
\]
es la proporción muestral de éxitos, y la ley de los grandes números afirma que
\[
\overline{X}_n \xrightarrow{P} p
\qquad
\text{e incluso}
\qquad
\overline{X}_n \xrightarrow{\mathrm{c.s.}} p.
\]

\subsection{Teorema Central del Límite}

La ley de los grandes números describe el comportamiento del promedio muestral \emph{sin escala}. El siguiente problema natural es entender las fluctuaciones de $\overline{X}_n$ alrededor de $\mu$ a una escala adecuada.

\begin{teo}[Teorema Central del Límite de Lindeberg--Lévy]
Si $X_1,X_2,\dots$ son i.i.d.\ con
\[
\E[X_1]=\mu,
\qquad
\Var(X_1)=\sigma^2\in(0,\infty),
\]
entonces
\[
\frac{\sum_{i=1}^n X_i - n\mu}{\sigma\sqrt{n}}
\xrightarrow{d}
\mathcal{N}(0,1).
\]
Equivalentemente,
\[
\frac{\sqrt{n}(\overline{X}_n-\mu)}{\sigma}
\xrightarrow{d}
\mathcal{N}(0,1).
\]
\end{teo}

\subsubsection*{Idea de la demostración mediante funciones características}

Definimos
\[
Y_i=\frac{X_i-\mu}{\sigma}.
\]
Entonces $Y_1,Y_2,\dots$ son i.i.d.\ y satisfacen
\[
\E[Y_1]=0,
\qquad
\Var(Y_1)=1.
\]

Sea
\[
S_n=\frac{1}{\sqrt{n}}\sum_{i=1}^n Y_i.
\]
Basta demostrar que
\[
S_n\xrightarrow{d}\mathcal{N}(0,1).
\]

Como las $Y_i$ son independientes,
\[
\varphi_{S_n}(t)
=
\E\left(e^{itS_n}\right)
=
\E\left(
e^{\,it\frac{1}{\sqrt{n}}\sum_{i=1}^n Y_i}
\right)
=
\E\left(
\prod_{i=1}^n e^{i\frac{t}{\sqrt{n}}Y_i}
\right).
\]

Usando independencia,
\[
\varphi_{S_n}(t)
=
\prod_{i=1}^n \E\left(e^{i\frac{t}{\sqrt{n}}Y_i}\right)
=
\left(\varphi_{Y_1}\left(\frac{t}{\sqrt{n}}\right)\right)^n.
\]

Ahora bien, como $\E[Y_1]=0$ y $\E[Y_1^2]=1$, se tiene el desarrollo alrededor de $0$:
\[
\varphi_{Y_1}(u)=1-\frac{u^2}{2}+o(u^2),
\qquad u\to 0.
\]

Tomando
\[
u=\frac{t}{\sqrt{n}},
\]
obtenemos
\[
\varphi_{Y_1}\left(\frac{t}{\sqrt{n}}\right)
=
1-\frac{t^2}{2n}+o\left(\frac{1}{n}\right).
\]

Por tanto,
\[
\varphi_{S_n}(t)
=
\left(1-\frac{t^2}{2n}+o\left(\frac{1}{n}\right)\right)^n
\to e^{-t^2/2}.
\]

Pero
\[
e^{-t^2/2}
\]
es la función característica de una variable $\mathcal{N}(0,1)$. Luego, por el Teorema de continuidad de Lévy,
\[
S_n\xrightarrow{d}\mathcal{N}(0,1).
\]

\subsubsection*{Interpretación}

El TCL afirma que, para $n$ grande,
\[
\overline{X}_n
\approx
\mathcal{N}\left(\mu,\frac{\sigma^2}{n}\right).
\]

Más precisamente, la cantidad
\[
\frac{\sqrt{n}(\overline{X}_n-\mu)}{\sigma}
\]
se comporta aproximadamente como una normal estándar cuando $n$ es grande.

Esta aproximación es una de las bases de la inferencia estadística clásica.

\subsection{Herramientas útiles para inferencia asintótica}

\subsubsection*{Teorema de mapeo continuo}

\begin{teo}[Teorema de mapeo continuo]
Si
\[
X_n\xrightarrow{d} X
\]
y $g:\mathbb{R}\to\mathbb{R}$ es continua, entonces
\[
g(X_n)\xrightarrow{d} g(X).
\]
\end{teo}

\medskip

\noindent
\textbf{Observación:} este resultado permite transportar convergencia en distribución a través de funciones continuas.

\subsubsection*{Teorema de Slutsky}

\begin{teo}[Teorema de Slutsky]
Si
\[
X_n\xrightarrow{d} X
\qquad\text{y}\qquad
Y_n\xrightarrow{P} c,
\]
donde $c$ es una constante, entonces:
\[
X_n+Y_n \xrightarrow{d} X+c,
\]
\[
X_nY_n \xrightarrow{d} cX,
\]
y, si $c\neq 0$,
\[
\frac{X_n}{Y_n}\xrightarrow{d}\frac{X}{c}.
\]
\end{teo}

\subsubsection*{Aplicación típica}

Supongamos que
\[
\frac{\sqrt{n}(\overline{X}_n-\mu)}{\sigma}
\xrightarrow{d}
\mathcal{N}(0,1),
\]
y además que $S_n$ es un estimador tal que
\[
S_n\xrightarrow{P}\sigma.
\]

Entonces, por Slutsky,
\[
\frac{\sqrt{n}(\overline{X}_n-\mu)}{S_n}
\xrightarrow{d}
\mathcal{N}(0,1).
\]

Este tipo de resultado aparece constantemente en inferencia estadística, cuando se reemplaza un parámetro desconocido por un estimador consistente.

\section{Copypaste de comandos}

\begin{verbatim}
\section{Título de la sección}             -> Crea una sección principal.

\subsection{Título de la subsección}       -> Crea una subsección dentro de una sección.

\textbf{texto en negrita}                  -> Pone el texto en negrita.

\emph{texto en énfasis}                    -> Da énfasis (normalmente cursiva).

\begin{defi}[Nombre de la definición]
  Texto de la definición.
\end{defi}
-> Entorno para escribir una definición.

\begin{teo}[Nombre del teorema]
  Texto del teorema.
\end{teo}
-> Entorno para enunciar un teorema.

\begin{obs}
  Texto de la observación.
\end{obs}
-> Entorno para agregar una observación o comentario.

\begin{enumerate}
  \item Primer ítem.
  \item Segundo ítem.
\end{enumerate}
-> Lista numerada; cada \item agrega un elemento.

$X_1, X_2, \ldots, X_n \iid f(\cdot\,;\theta)$   -> Notación de muestra i.i.d. con distribución f.

$\E[X]$                                        -> Valor esperado de X.

$\bar{X}_n$                                    -> Media muestral.

$\hat{\theta}$                                -> Estimador de theta.

$\sigma^2$                                     -> Varianza poblacional.

$\theta \in \Theta$                          -> El parámetro theta pertenece al espacio paramétrico Theta.

\[
  \bar{X}_n = \frac{1}{n}\sum_{i=1}^{n} X_i
\]
-> Fórmula de la media muestral usando suma.

\[
  \bar{X}_n \xrightarrow{\text{c.s.}} \mu
\]
-> Convergencia casi segura de la media muestral hacia mu.
\end{verbatim}
